\section{Architecture}

\subsection{CUA Abstraction: A Closed-Loop Perception--Brain--Action System}

Modern Computer-Use Agents (CUAs) have evolved from early rule-based
heuristics into autonomous systems driven by multimodal large language models
(MLLMs). Synthesizing recent core studies, the foundational CUA framework can
be abstracted as a three-layer \emph{Perception--Brain--Action} (PBA)
architecture \cite{wang2024llm_agents_survey}. This architecture operates
across both GUI-centric and API-centric interaction pathways, while increasingly
coupling security mechanisms into the execution loop. The unified framework is
illustrated in Fig.~\ref{fig:cua_architecture}.

\begin{figure*}[t]
    \centering
    \includegraphics[width=\textwidth]{images/picture3.1.png}
    \caption{Three-stage CUA architecture under a closed-loop
    Perception--Brain--Action design. Stage 1 performs perception and
    sanitization, Stage 2 handles planning and reflection, and Stage 3 executes
    validated actions under sandboxed isolation.}
    \label{fig:cua_architecture}
\end{figure*}

\subsubsection{Perception: Alignment and State Representation in Heterogeneous Environments}

The perception layer transforms high-dimensional and noisy digital inputs (e.g.,
pixel streams, HTML sources, API responses) into structured state
representations consumable by foundation models. Technically, this layer has
shifted from metadata-dependent parsing to vision-native grounding.

\textbf{Structured text and DOM-based perception.}
Early web agents relied heavily on underlying interface metadata. WebArena and
Mind2Web used reduced DOM trees or accessibility trees as primary inputs
\cite{zhou2024webarena,deng2024mind2web}. While efficient in instrumented web
settings, such approaches are brittle in closed-source software and dynamic
rendering scenarios. In machine-centric interaction settings, Gorilla further
examined structured API-document and JSON-schema grounding for tool-use
decisions \cite{patil2023gorilla}.

\textbf{Pure visual grounding and multimodal localization.}
To remove dependence on privileged interface metadata, recent work moved toward
direct screenshot understanding. CogAgent and ScreenAI demonstrated that
high-resolution vision-language modeling can capture complex UI layout and
semantics \cite{hong2024cogagent,baechler2024screenai}. SeeClick and Ferret-UI
then introduced coordinate-grounded objectives, significantly improving
fine-grained element localization on desktop and mobile environments
\cite{cheng2024seeclick,you2024ferretui}. More recently, OmniParser proposed a
pure-vision parsing paradigm that tokenizes UI screenshots and extracts icon
semantics and text bounding boxes without external OCR dependency
\cite{lu2024omniparser}.

\textbf{Security-aware perception and environment sanitization.}
As CUAs are exposed to adversarial environments, perception-layer fragility has
become explicit. AgentXploit showed that hidden or context-poisoned interface
content can trigger indirect prompt injection and manipulate downstream agent
behavior \cite{liu2024agentxploit}. To mitigate this risk, Chen et al.
introduced visual pre-filtering and sanitization pipelines for VLM-based CUAs,
isolating untrusted external signals before they enter planning
\cite{chen2024defending_ipi_vlm}.

\subsubsection{Brain: Long-Horizon Planning, Memory, and Reflection}

The brain layer serves as the system's executive control module: it handles
goal decomposition, intent grounding, and dynamic routing across heterogeneous
action channels. Its central challenge is maintaining coherent reasoning under
long-horizon execution, where catastrophic forgetting and objective drift are
common.

\textbf{Temporal planning and dynamic routing.}
Building on the ReAct paradigm of interleaving reasoning and acting
\cite{yao2023react}, recent CUA systems adopt more explicit planning topology.
OSCopilot presented a general operating-system agent framework that decomposes
system-level objectives into executable subtasks \cite{wu2024oscopilot}. In
hybrid tool/UI settings, the brain must also perform cost-aware routing:
prefer direct API/tool invocations when available (Toolformer-style), and
degrade gracefully to GUI interaction when no reliable interface exists
\cite{schick2023toolformer}.

\textbf{Episodic memory and exploration.}
When operating in unfamiliar software environments, CUAs benefit from active
exploration and memory accumulation. AppAgent introduced autonomous trial-and-
error interaction to build episodic action--state memory for downstream reuse
\cite{zhan2024appagent}. AgentBoard evaluations further indicate that external
memory mechanisms (e.g., retrieval over trajectory traces) can improve success
rates in long-context, multi-turn tasks \cite{ma2024agentboard}.

\textbf{Closed-loop reflection and alignment guardrails.}
Real-world execution includes latency, UI non-determinism, and system errors.
Reflexion introduced verbal reinforcement and self-correction loops that allow
agents to backtrack and revise plans after failed actions \cite{shinn2023reflexion}.
To reduce long-horizon goal hijacking risk, recent secure-agent designs also
emphasize guardrails-by-construction, where side-channel monitors continuously
check whether decision trajectories remain aligned with user intent
\cite{liu2024agentxploit}.

\subsubsection{Action: Unified Action Space and Polyglot Sandbox Isolation}

The action layer compiles planned intent into concrete operational effects. Its
core design tension is balancing cross-platform generalization against
potentially destructive execution risk.

\textbf{End-to-end native GUI control.}
Traditional CUA stacks often depended on manually engineered controllers for
coordinate mapping and event translation. UI-TARS and follow-up systems
promoted an end-to-end paradigm that directly outputs absolute/relative
coordinates and keyboard/mouse events, reporting strong benchmark performance
\cite{qin2025uitars}. OS-ATLAS further showed that large-scale cross-platform
pretraining (Windows/macOS/Linux) can learn a more unified and generalizable
action space \cite{boyce2025osatlas}.

\textbf{Code-as-Action and Agent-Computer Interface (ACI).}
For backend-heavy or computation-intensive tasks, pure GUI clicking is often
inefficient. SWE-agent proposed an agent-computer interface (ACI) that allows
models to generate and execute Bash/Python code in terminal environments,
expanding CUA capability in software engineering workflows
\cite{yang2024sweagent}.

\textbf{Isolated execution and human-in-the-loop (HITL).}
Because action-layer outputs can directly mutate file systems, credentials, and
network states, environment-level isolation is critical. OSWorld and related
evaluation settings highlight the necessity of realistic yet controlled
execution environments \cite{xie2024osworld}. In production-grade deployments,
generated code should run in stateless containers or virtual machines; in
addition, sensitive operations (e.g., privilege escalation) should trigger
human-in-the-loop confirmation to provide strong validation for irreversible
actions.
